% ch7_b §7.3–7.5 (书页 219–227 / p233–p241)
%--A-TAIL--
%JOIN%其中 $m_{e}$ 是电子的有效质量——空穴亦然。这些公式表明，迁移率可以是温度的函数；随着 $T$ 升高，分布展宽，平均弛豫时间随之改变。在金属中，$\tau$ 随能量的变化关系影响不大；重要的只是 $\tau(\mathcal{E}_{F})$ 的数值。

\section{弛豫时间的计算}

我们至今还没有求解积分方程 (7.16)。事实上，它最普遍的可能解应当写成
\begin{equation*}
g_{\mathbf{k}}=\left(-\frac{\partial f^{0}}{\partial\mathcal{E}}\right)e\mathbf{E}\cdot\boldsymbol{\Lambda}(\mathbf{k}),\tag{7.40}
\end{equation*}
其中 $\boldsymbol{\Lambda}(\mathbf{k})$ 是在 费米面上每一点 $\mathbf{k}$ 处定义的一个矢量。我们的初等解 (7.19) 相当于取
\begin{equation*}
\boldsymbol{\Lambda}(\mathbf{k})=\tau\,\mathbf{v}_{\mathbf{k}},\tag{7.41}
\end{equation*}
这表明 $\boldsymbol{\Lambda}$ 是电子的\emph{矢量平均自由程}（vector mean free path）。在一般情形下，$\boldsymbol{\Lambda}(\mathbf{k})$ 的大小可以变化，而且在 费米面上还可以偏离 $\mathbf{v}_{\mathbf{k}}$ 的方向。

有时人们假定可以写成
\begin{equation*}
\boldsymbol{\Lambda}(\mathbf{k})=\tau(\mathbf{k})\,\mathbf{v}_{\mathbf{k}},\tag{7.42}
\end{equation*}
其中 $\tau(\mathbf{k})$ 是一个在 费米面上变化的\emph{各向异性弛豫时间}（anisotropic relaxation time）。然而容易证明，这并不总是积分方程的完整解——而且也不存在计算函数 $\tau(\mathbf{k})$ 的直接方法。

事实上，唯一简单的解就是初等解 (7.19)。假设我们把它作为 $g_{\mathbf{k}}$ 代入 (7.16)，并且还假定散射是弹性的，即
\begin{equation*}
Q(\mathbf{k},\mathbf{k}')\,\mathrm{d}\mathbf{k}'=\delta(\mathcal{E}-\mathcal{E}')\,\mathcal{Q}(\mathbf{k},\mathbf{k}')\,\mathrm{d}\Omega'\,\mathrm{d}E,\tag{7.43}
\end{equation*}
其中 $\mathrm{d}\Omega'$ 是散射后 $\mathbf{k}'$ 的方向上的立体角元（$\mathbf{k}'$ 的大小现在由它必须与 $\mathbf{k}$ 具有相同能量 $\mathcal{E}'$ 这一要求所确定）。

在两边消去能量的 delta 函数之后，我们得到
\begin{equation*}
\mathbf{v}_{\mathbf{k}}\cdot\mathbf{E}=\tau\int(\mathbf{v}_{\mathbf{k}}-\mathbf{v}_{\mathbf{k}'})\cdot\mathbf{E}\,\mathcal{Q}(\mathbf{k},\mathbf{k}')\,\mathrm{d}\Omega'\tag{7.44}
\end{equation*}
积分在 费米面上进行。这是一个泛函关系，它对 $\mathcal{Q}(\mathbf{k},\mathbf{k}')$ 的形式施加了若干条件。容易证明：当我们有球形 费米面、$|\mathbf{v}_{\mathbf{k}}|$ 为常数，并且
\begin{equation*}
\mathcal{Q}(\mathbf{k},\mathbf{k}')=\mathcal{Q}(\theta);\tag{7.45}
\end{equation*}
也就是说，当散射概率仅是两个波矢之间夹角的函数时，关系式 (7.44) 便能够成立。这个证明只不过是球面三角学的一道练习题，我们把它留给读者。

如果这些条件得到满足，我们立刻得到
\begin{equation*}
\frac{1}{\tau}=\int(1-\cos\theta)\,\mathcal{Q}(\theta)\,\mathrm{d}\Omega'.\tag{7.46}
\end{equation*}
这表明，弛豫时间反比于一个对所有过程的散射概率的积分——但带上了因子 $(1-\cos\theta)$ 的权重，向大角散射倾斜。正如在 (7.44) 中所见，这个因子来自 $(\mathbf{v}_{\mathbf{k}}-\mathbf{v}_{\mathbf{k}'})\cdot\mathbf{E}$；重要的并不是电子被散射这件事，而是在此过程中电子速度沿电场方向的分量改变了多少。

我们可以借助数密度为 $N_{i}$ 的杂质的微分散射截面 $\sigma(\theta)$ 来表达散射概率：平均自由程 $\Lambda$ 由下式给出
\begin{equation*}
\frac{1}{\Lambda}=N_{i}\,2\pi\int_{0}^{\pi}(1-\cos\theta)\,\sigma(\theta)\sin\theta\,\mathrm{d}\theta.\tag{7.47}
\end{equation*}

\section{杂质散射}

现在，有了 (6.70)、(7.25) 和 (7.47)，我们就具备了从第一性原理出发计算含带电杂质的自由电子金属之电导率所需的材料。我们通常把它表示成\emph{电阻率}
\begin{align*}
\rho&=\frac{mv_{F}}{ne^{2}\Lambda}\\
&=\frac{mv_{F}}{ne^{2}}\,N_{i}\,2\pi\int_{0}^{\pi}(1-\cos\theta)\left(\frac{2mZe^{2}}{\hbar^{2}}\right)^{2}\frac{\sin\theta\,\mathrm{d}\theta}{(K^{2}+\lambda^{2})^{2}}\\
&=\frac{mv_{F}}{ne^{2}}\,N_{i}\,2\pi\left(\frac{2mZe^{2}}{\hbar^{2}\lambda^{2}}\right)^{2}\int_{0}^{1}\frac{8z^{3}\,\mathrm{d}z}{\left\{1+(2k_{F}/\lambda)^{2}z^{2}\right\}^{2}}.\tag{7.48}
\end{align*}
这里我们利用了散射的几何关系，即 $K=2k_{F}\sin\frac{1}{2}\theta$，并在被积函数中令 $z=\sin\frac{1}{2}\theta$。

(7.48) 中的积分作为 $(2k_{F}/\lambda)$ 的函数很容易算出——我们不打算写出精确的公式。要注意的要点是：一个带电杂质的行为就像一个半径为
\begin{equation*}
R\sim\left(\frac{2mZe^{2}}{\hbar^{2}\lambda^{2}}\right)\tag{7.49}
\end{equation*}
的几何障碍物；由 (5.22) 和 (3.6)，它与原子球的半径相当。我们还注意到，这一电阻与温度无关，而且应当按价差 $Z$ 的平方变化（\emph{Linde 定则}，Linde's rule）。

实际上，由 (7.48) 给出的电阻太大了。更好的做法是采用分波公式 (5.40) 来计算散射截面。对 $\theta$ 的积分可以做出来；我们得到
\begin{equation*}
\frac{1}{\Lambda}=N_{i}\,\frac{4\pi}{k_{F}^{2}}\sum_{l=1}^{\infty}l\sin^{2}(\eta_{l-1}-\eta_{l}).\tag{7.50}
\end{equation*}
然后可以确定各相移，使之满足 Friedel 求和规则 (5.43)。任何在 费米能级附近具有 $d$ 共振（§§5.5、6.10）的过渡金属杂质，显然必定产生很大的\emph{剩余电阻率}（residual resistivity），此时 $\eta_{2}\sim\pi/2$。对磁性杂质而言，依赖于自旋的 s--d 相互作用还会导致\emph{电阻极小}（resistance minimum），即 \emph{Kondo 效应}（Kondo effect），这将在 §10.6 中讨论。

要计算半导体中与带电杂质散射相联系的迁移率，我们可以采用与 (7.48) 本质上相同的论证。此时弛豫时间实际上是依赖于能量的：
\begin{equation*}
\frac{1}{\tau(\mathcal{E})}=N_{i}\,v\cdot 2\pi\left(\frac{Ze^{2}}{\epsilon\mathcal{E}}\right)^{2}\int_{0}^{1}\frac{\frac{1}{2}z^{3}\,\mathrm{d}z}{\left\{z^{2}+\lambda^{2}\hbar^{2}/8m^{*}\mathcal{E}\right\}^{2}},\tag{7.51}
\end{equation*}
其中 $\lambda$——现在由 (5.26) 给出——被假定远小于热激发所容许的 $k$ 值范围。

我们若把这个积分算出来，就会发现它随 $\mathcal{E}$ 只变化得很慢：$\tau$ 通过因子 $\mathcal{E}^{2}/v\propto\mathcal{E}^{\frac{3}{2}}m^{*\frac{1}{2}}$ 依赖于载流子的能量。当我们把它代入 (7.38) 时，实际上发现这相当于用平均值 $kT$ 代替 $\mathcal{E}$。于是，对于杂质散射，
\begin{equation*}
\mu\propto T^{\frac{3}{2}}m^{*\frac{3}{2}}.\tag{7.52}
\end{equation*}
（译注：按式 (7.51) 的 $\tau\propto\mathcal{E}^{3/2}m^{*1/2}$ 代入 (7.38)、(7.39)，
这类公式确实非常粗糙，但“迁移率应当随温度升高而增大”这一普遍的定性论证是可靠的。我们的意思是：载流子气越热，平均载流子就运动得越快，也就越不容易被带电杂质散射。

\section{“理想”电阻}

计算金属的所谓\emph{“理想”电阻}（'ideal' resistance）——即来源于晶格热振动之散射的电阻——是一个更为复杂的问题。我们只能把论证的主要线索勾画出来。

作为一级近似，我们可以利用 §6.14 中引入的形变势原理。电子感受到一个涨落势，其大小为
\begin{equation*}
\delta\mathcal{E}=\tfrac{2}{3}\mathcal{E}_{F}\,\Delta,\tag{7.53}
\end{equation*}
其中 $\Delta$ 是局域的、随热涨落的体膨胀。因此，电阻应当正比于密度的均方涨落；按照经典统计力学的一个熟知原理，它由下式给出
\begin{equation*}
\overline{\Delta^{2}}=NkT\beta,\tag{7.54}
\end{equation*}
其中 $\beta$ 是压缩率。

如果用声速 $s$ 和密度 $D$ 来表示，就有
\begin{equation*}
\overline{\Delta^{2}}=\frac{NkT}{Ds^{2}}=\frac{kT}{Ms^{2}},\tag{7.55}
\end{equation*}
其中 $M$ 是一个离子的质量。这个结果更为人们熟悉的形式是如下的一般关系
\begin{equation*}
\rho_{t}\propto\frac{T}{M\Theta^{2}}\tag{7.56}
\end{equation*}
它表明电阻正比于绝对温度。电阻随原子质量以及随 Debye 温度 $\Theta$（如 §2.4 中那样，$\Theta$ 与 $s$ 成正比）的变化，大体上与实验相符。

为了在这一结果上有所改进，我们可以采用 (6.87)，即晶格中每个离子的有效微分截面。由 (7.47) 我们得到平均自由程
\begin{equation*}
\frac{1}{\Lambda_{t}}=N\,2\pi\int_{0}^{\pi}\sigma_{a}(\theta)\,N|\mathbf{K}\cdot\mathbf{U}_{\mathbf{q}}|^{2}(1-\cos\theta)\sin\theta\,\mathrm{d}\theta,\tag{7.57}
\end{equation*}
其中 $\sigma_{a}(\theta)$ 是“自由原子”的微分散射截面。

假设我们只考虑 N 过程，即散射矢量 $\mathbf{K}$ 等于声子波矢 $\mathbf{q}$ 的过程。在 (2.109) 与 (2.110) 中已经算出因子 $|\mathbf{K}\cdot\mathbf{U}_{\mathbf{q}}|^{2}$。在高温下
\begin{align*}
N|\mathbf{K}\cdot\mathbf{U}_{\mathbf{q}}|^{2}&\approx\frac{kTK^{2}}{M\nu_{\mathbf{q}}^{2}}\\
&\approx\frac{kT}{Ms^{2}}\\
&\approx\frac{\hbar^{2}q_{D}^{2}\,kT}{Mk^{2}\Theta^{2}}\tag{7.58}
\end{align*}
——以 Debye 温度表出。实际上，这就是 (7.54) 与 (7.55) 中的 $\overline{\Delta^{2}}$。

由于 (7.58) 与 $\mathbf{K}$ 无关，也就是与散射角 $\theta$ 无关，我们可以把这个积分单独算出来，并写成
\begin{equation*}
\frac{1}{\Lambda_{t}}=N\bar{\sigma}_{a}\,\frac{\hbar^{2}q_{D}^{2}\,kT}{Mk^{2}\Theta^{2}},\tag{7.59}
\end{equation*}
其中
\begin{equation*}
\bar{\sigma}_{a}\equiv 2\pi\int\sigma_{a}(\theta)\,(1-\cos\theta)\sin\theta\,\mathrm{d}\theta,\tag{7.60}
\end{equation*}
即孤立原子的总散射截面，并以因子 $(1-\cos\theta)$ 计及沿场方向动量的变化。

利用 Lindemann 熔化公式 (2.117)，这个公式还可以进一步约化。我们得到
\begin{equation*}
\frac{1}{\Lambda_{t}}\approx N\bar{\sigma}_{a}\cdot\tfrac{2}{3}x_{m}^{2}\,\frac{T}{T_{m}},\tag{7.61}
\end{equation*}
其中 $x_{m}$ 是 (2.117) 中量级为 0.2 的常数，$T_{m}$ 是熔化温度。我们还可以再往前走一点，假定一个离子的散射截面与带电杂质的散射截面大体相同。对 (7.49) 稍作变换，就得到大致如下的结果
\begin{equation*}
\Lambda_{t}\sim 50a\,\frac{T_{m}}{T},\tag{7.62}
\end{equation*}
其中 $a^{3}$ 是一个晶胞的体积；金属中电子的平均自由程在熔点处应当具有晶格常数 50 倍的量级。这并不严格正确——但作为数量级的估计还是站得住脚的。

要得到电阻本身，我们应当把 (7.59) 或 (7.62) 与 (7.33) 结合起来，写成
\begin{equation*}
\rho_{t}=\frac{mv_{F}}{ne^{2}\Lambda_{t}},\tag{7.63}
\end{equation*}
这对自由电子气近似地成立。这些公式使我们粗糙的函数关系 (7.56) 有了具体的数字。

但在低温下情况要更复杂。(7.58) 这个结构因子是在经典统计的假设下算出来的。如果进入低温区，我们就必须把比如 $(kT/\hbar\nu_{\mathbf{q}})$ 换成由 (2.46) 给出的 $\bar{n}_{\mathbf{q}}$。\footnote{零点能并不进入这个公式；要证明这一点，需要对散射作细致的分析，把声子的发射过程与吸收过程区分开来。}我们得到
\begin{equation*}
N|\mathbf{K}\cdot\mathbf{U}_{\mathbf{q}}|^{2}\approx\frac{kT}{Ms^{2}}\left\{\frac{\hbar\nu_{\mathbf{q}}/kT}{\exp(\hbar\nu_{\mathbf{q}}/kT)-1}\right\}.\tag{7.64}
\end{equation*}

(7.57) 积分中的这个因子起到了在角度变量上设置切断的作用。当我们需要吸收或发射频率高得无法被热激发的声子时，散射概率迅速下降；也就是说，对于满足
\begin{equation*}
\hbar\nu_{\mathbf{q}}>kT.\tag{7.65}
\end{equation*}
的散射，我们实际上不予考虑。

对于 Debye 模型以及球形 费米面上的 N 过程，这种情况发生于如下的 $q$ 值和散射角 $\theta$，使得
\begin{equation*}
\frac{q}{2k_{F}}=\sin\tfrac{1}{2}\theta>\frac{q_{D}}{2k_{F}}\,\frac{T}{\Theta}.\tag{7.66}
\end{equation*}

如果我们假定 $\sigma_{a}(\theta)$ 并非随 $\theta$ 急剧变化的函数，因而可以作为一个常数提到积分之外，就能相当精确地计算平均自由程。于是 (7.57) 给出
\begin{equation*}
\frac{1}{\Lambda_{t}}=N\bar{\sigma}_{a}\,\frac{\hbar^{2}q_{D}^{2}\,kT}{Mk^{2}\Theta^{2}}\left(\frac{T}{\Theta}\right)^{4}\int_{0}^{\Theta/T}\frac{4z^{4}\,\mathrm{d}z}{(e^{z}-1)},\tag{7.67}
\end{equation*}
其中我们已令
\begin{equation*}
z=\frac{\hbar\nu_{\mathbf{q}}}{kT}=\left(\frac{q}{q_{D}}\right)\left(\frac{\Theta}{T}\right).\tag{7.68}
\end{equation*}

%FIG{123}: {电子--声子 N 过程：（a）高温下；（b）低温下。}

在高温下，$\Theta/T$ 很小，积分趋于 $(\Theta/T)^{4}$，我们就回到了 (7.59)。在低温下，积分趋于一个常数 124.4，于是电阻正比于 $T^{5}$。“理想”电阻的这种强烈的温度依赖性是一种典型的量子效应，可与 Debye $T^{3}$ 比热定律相比拟。

我们可以看出发生了什么。有效散射截面是一个对角度的积分，带有因子
\begin{equation*}
(1-\cos\theta)\sin\theta\,\mathrm{d}\theta=8\sin^{3}\tfrac{1}{2}\theta\,\mathrm{d}(\sin\tfrac{1}{2}\theta)\tag{7.69}
\end{equation*}
因而对 (7.64) 所提示的切断非常敏感。如果取 (7.66) 作为切断角，我们就应当发现一个正比于
\begin{equation*}
\int_{0}^{(q_{D}/2k_{F})(T/\Theta)}8\sin^{3}(\tfrac{1}{2}\theta)\,\mathrm{d}(\sin\tfrac{1}{2}\theta)\propto\left(\frac{T}{\Theta}\right)^{4},\tag{7.70}
\end{equation*}
的因子，这正是 (7.67) 中出现的那一个。长波声子作为电阻的来源效果很差，因为它们只能把电子散射过一个小角——而这个角随温度成比例地减小。

把 (7.67) 与 (7.63) 结合起来得到的公式，就是所谓的 \emph{Bloch--Grüneisen 定律}（Bloch--Grüneisen law）。许多金属在低温下近似遵从它——但不应对它过分拘泥于字面，理由如下：

（i）我们假定了电子被弹性散射。这并不真实——但如果我们按照电子--声子的产生与湮没过程作一个形式上的分析，并计入电子失去或获得的能量，就得到几乎相同的结果。(7.67) 中的积分变为
\begin{equation*}
\int_{0}^{\Theta/T}\frac{4z^{5}\,\mathrm{d}z}{(e^{z}-1)(1-e^{-z})},
\end{equation*}
它在 $T$ 很小和很大时的渐近行为都与 (7.67) 相同。

（ii）Debye 谱的假设只是一种方便的过分简化。纵模与横模各自的贡献应当分别考虑。

（iii）函数 $\sigma_{a}(\theta)$ 不是常数，因此不能把它从 (7.57) 的积分号下提出来。计算 $\sigma_{a}(\theta)$ 的公式由 (6.89) 与 (6.93) 给出。特别是，当 $\theta$ 超过大约 $20^{\circ}$--$30^{\circ}$ 以后，$\sigma_{a}(\theta)$ 趋于减小——这也对电阻随温度的变化产生影响。

（iv）最严重的误差在于忽略了 U 过程。如 §2.7 所示，如果 $\mathbf{K}$ 越出了某个 Brillouin 区的边界，我们就应当写成
\begin{equation*}
\mathbf{q}=\mathbf{K}-\mathbf{g},\tag{7.71}
\end{equation*}
其中 $\mathbf{g}$ 是一个倒格矢。在我们的公式 (7.57) 中，并没有什么特殊理由使这样的 $\mathbf{K}$ 值无法达到；原子的散射截面作为角度的函数会一直延伸到 $\theta=\pi$，那将使 $K=2k_{F}$，大于 Debye 极限 $q_{D}$。

U 过程给电阻造成的困难只不过在于，它们会引起复杂的几何问题。即使在高温下，如果把 (7.71) 代入 (7.58)，我们也会发现“结构因子”不再与散射角无关：它包含一个额外的因子 $K^{2}/q^{2}$，这个因子可以远大于 1。我们可以在图 124(a) 中看到这一点，其中 $K$ 显然远大于 $q$。对于这样大的 $\theta$ 值，微分散射截面 $\sigma_{a}(\theta)$ 也许很小——但对电阻的贡献却被因子 $(1-\cos\theta)$ 重重地加权。

某些 U 过程所需要的很小的 $q$ 值，在我们考虑电阻随温度的变化时也很重要。正如 (7.66) 所表明的，直到温度降到很低，这些过程才被“冻结”。如果把图 124(a) 的矢量三角形在重复区图式中重新画出，这一点可以看得更清楚。这时，一个倒逆过程所要求的 $q$ 的最小值，显然就是一个区中的 费米面与其在相邻区中的重复像之间的最小距离。

%FIG{124}: {电子--声子 U 过程：（a）在自由电子球上；（b）在重复区图式中。}

如果 费米面朝区边界的方向鼓出，那么 $q$ 的最小值就减小——于是 (7.57) 中带因子 $K^{2}/q^{2}$ 的贡献将得到增强。的确，看起来几乎就好像：倘若 费米面碰到区边界，$\rho_{i}$ 中就会出现一个奇点——不过，正如 §6.13 中指出的，对于这样的过程，电子--声子相互作用的矩阵元自动为零。

U 过程对电阻的贡献没有简单的公式。不过，给出 (7.56) 和 (7.67) 的那些同样的普遍因素仍在起作用——尽管后一个表达式作为精确的理论结果并不成立。我们能够理解这一现象的全部物理特征，但必须依靠繁冗的数值计算才能得到答案。举一点为例：按照 (6.93)，原子在 $\mathbf{K}$ 值大（即 $\theta$ 大）时的有效截面依赖于 $\mathcal{V}'$ 的大小，即赝原子的“芯赝势”（'core pseudo-potential'）。这应当反映在电阻的大小上——但电阻对 $\mathcal{V}'$ 的敏感，既可能来自它使 费米面畸变、使区边界附近的波函数相混合的方式，也同样可能来自它对散射矩阵元的直接贡献。

%FIG{125}: {液态金属的典型结构因子，图中标出不同价数下的 $2k_{F}$，以及温度依赖性。}

对\emph{液态}金属的类似计算实际上要容易得多。假定一个 N.F.E. 模型是合理的，此时 (6.88)、(6.91)、(7.25) 和 (7.47) 引导我们得到公式
\begin{equation*}
\rho_{L}=\frac{3\pi}{4}\,\frac{1}{\hbar^{2}e^{2}v_{F}^{2}k_{F}^{4}}\,\frac{1}{N}\int_{0}^{2k_{F}}|w_{s}(K)|^{2}\,S(K)\,K^{3}\,\mathrm{d}K.\tag{7.72}
\end{equation*}
在这个表达式中，起实质性作用的要素是赝原子形状因子 $|w_{s}(K)|^{2}$ 和液体结构因子 $S(K)$，后者可以用中子或 X 射线衍射在实验上测定。事实上，大多数液态金属中原子的排列方式如此相似，以致在按原子半径标度之后，$S(K)$ 可以近似地用一条统一的标准曲线来表示（图 125）。在小的 $K$ 处，结构因子很小，因为原子分布在密度上近乎均匀——事实上，经典公式 (7.54) 在这里仍然成立。但第一
